%% ch07_results.tex  —  Chapter 7: Experimental Results
%% HSP-Agile Thesis, East China Normal University, 2026

\chapter{Experimental Results}
\label{ch:results}

This chapter is evidence for the three Tier-1 claims
(Table~\ref{tab:claim-map}), not method restatement.
Mechanism first (E3--E6: where and why the ideas work), then aggregates
(E1/E2), then cost and transfer.
Table~\ref{tab:results-roadmap} is the map; unless noted, $n{=}120$ hard tasks
and the ten compared modes apply.

\begin{inlinetable}{Results roadmap: mechanism-first presentation order.}
  \label{tab:results-roadmap}
  \footnotesize
  \begin{tabular}{lll}
    \toprule
    RQ & Section & Key figure / table \\
    \midrule
    RQ4 & \ref{sec:results:rq4} & Fig.~\ref{fig:complexity} (E3) \\
    RQ3 & \ref{sec:results:rq3} & Figs.~\ref{fig:convergence},~\ref{fig:feedback-variant} \\
    RQ1 & \ref{sec:results:rq1} & Table~\ref{tab:main-results} \\
    RQ2 & \ref{sec:results:rq2} & Figs.~\ref{fig:prevention-heatmap},~\ref{fig:prevention-bar} \\
    RQ5 & \ref{sec:results:rq5} & Fig.~\ref{fig:pareto} \\
    \bottomrule
  \end{tabular}
\end{inlinetable}

%=============================================================================
\section{RQ4: Specification Complexity Scaling (E3)}
\label{sec:results:rq4}
%=============================================================================

\textbf{Claim (C1).}
If prevention depends on semantic regions, the M--baseline gap should grow
with guard-overlap density---where random tests fail.

\begin{figure}[t]
  \centering
  \widefig{figures/complexity_heatmap.pdf}
  \caption{\textbf{Claim:} M's advantage over B1/B2 grows with overlap density.
    \textbf{How to read:} rows are density tertiles; cells are mean strict
    formal conformance.  Compare columns within the high-overlap row.}
  \label{fig:complexity}
\end{figure}

Where does the advantage live?
In the sparse-overlap tier, M/B2/B1 are close ($\approx$93/91/88\%).
In the dense-overlap tier, M leads B2 by $\approx$3\,pp and B1 by $\approx$8\,pp
(87 vs.\ 84 vs.\ 79\%).
The gap widens exactly where random tests miss boundary regions---the same
failure mode as \code{classify\_signal}'s $(-3,-10)$, now nested as
$a>80\land b<20$ before $a>50$.
Under B1 the broader guard wins; under M the Semantic Feedback IR names S1 and
the second attempt restores order.

\paragraph{Answer to RQ4 / C1.}
The advantage scales with overlap density ($\geq$3\,pp over B2, $\geq$8\,pp
over B1 in the dense tier): semantic-region witnesses matter most where
sampling fails.

%=============================================================================
\section{RQ3: Component Attribution (E5, E6, Ablation)}
\label{sec:results:rq3}
%=============================================================================

\subsection{Repair Dynamics (E5)}

If most of the conformance gain arrives by the second attempt, budget $K=3$
is a sweet spot; if not, the extra call is wasted money.
Figure~\ref{fig:convergence} settles that question.

\begin{figure}[t]
  \centering
  \widefig{figures/repair_convergence.pdf}
  \caption{\textbf{Claim:} most repair gain arrives by iteration~2, justifying
    budget $K=3$.
    \textbf{How to read:} solid line is mode~M; compare to B2 at the same
    attempt index.  Shaded bands are 95\% bootstrap CIs.}
  \label{fig:convergence}
\end{figure}

The $\mathrm{Conf}(k)$ trajectory, extracted from 258 backfilled M/B2
\texttt{attempt\_history} records, settles it: mode~M rises from 78.3\% at
$k=1$ to 84.5\% at $k=2$, then plateaus at 83.4\% at $k=3$.
Most of the repair gain arrives within two iterations; further calls show
diminishing returns on this benchmark.

\subsection{Feedback Variant Isolation (E6)}

\begin{figure}[t]
  \centering
  \widefig{figures/feedback_variant_bar.pdf}
  \caption{\textbf{Claim:} specification-level structure, not prompt volume, drives
    repair.
    \textbf{How to read:} A~=~test only (no scenario info); B~=~test+expected;
    C~=~full Semantic Feedback IR.  Compare bar heights, not bar widths.}
  \label{fig:feedback-variant}
\end{figure}

\paragraph{Vignette (feedback content).}
On the \code{classify\_signal}-style witness $(-3,-10)$, variant~A tells the
model only that the test failed; variant~B adds the expected string
\texttt{Error}; variant~C additionally names Scenario~S1 and the guard
$\mathit{level}<0$.
The +7.7\,pp gap from A to C is the empirical cost of omitting that structure.

Figure~\ref{fig:feedback-variant} isolates feedback content under otherwise
identical mode~M settings ($n=120$ tasks per variant; Table~\ref{tab:feedback-e6}).
Variant~C (full Semantic Feedback IR) reaches mean Conf\,=\,86.9\%---a gain of
\textbf{+7.7\,pp} over test-only feedback (A, 79.1\%) and \textbf{+8.0\,pp} over
test-plus-expected (B, 78.9\%).
A and B are statistically indistinguishable ($\Delta < 0.3$\,pp).
Adding the expected string alone is not a substitute for naming the violated
scenario and guard.
The repair benefit attributed to mode~M in E5 is therefore structural, not a
side effect of longer prompts.

\begin{inlinetable}{E6 feedback variant summary, $n=120$ per variant.}
  \label{tab:feedback-e6}
  \begin{tabular}{lccc}
    \toprule
    Variant & Feedback surface & Mean Conf & Mean iterations \\
    \midrule
    A & test only           & 79.1\% & 2.98 \\
    B & test + expected     & 78.9\% & 2.97 \\
    C & full Semantic IR    & 86.9\% & 2.99 \\
    \bottomrule
  \end{tabular}
\end{inlinetable}

\subsection{Ablation Study (E1/RQ3)}

The ablation study removes one HSP-Agile component at a time, holding all other
parameters fixed, to isolate each component's contribution to the overall formal
conformance of 90.4\%.
We define the ablation delta for component $x$ as
$\Delta_x = \overline{\mathit{Conf}}_{\text{ablated}} - \overline{\mathit{Conf}}_M$;
a negative delta indicates that removing the component \emph{hurts} conformance.

\begin{figure}[t]
  \centering
  \widefig{figures/ablation_contribution_ci.pdf}
  \caption{\textbf{Key observation:} repair earns the largest conformance
    share (A3 $-$6.3\,pp); the checker is next (A1 $-$3.9\,pp); the pattern
    guard is a prevention gate, not a Conf.\ booster.
    \textbf{How to read:} bars are $\Delta$ vs.\ full M with 95\% bootstrap CIs.}
  \label{fig:ablation}
\end{figure}

Which stage earns the gain?
A3 (no repair) $-$6.3\,pp; A1 (no checker) $-$3.9\,pp; A2 (no pattern guard)
$+$0.5\,pp on Conf.\ but lower structural PDR---prevention value that
conformance alone misses (Section~\ref{sec:results:rq2}).

\subsection{Ranking of Component Contributions}

Repair $>$ checker $\gg$ pattern guard on Conf.\ (A3/A1/A2);
the guard's value shows up in structural PDR, not in mean Conf.

\paragraph{Answer to RQ3 / C2--C3.}
Typed feedback is the largest repair lever (+7.7\,pp, E6; A3 $-$6.3\,pp);
the checker is next (A1 $-$3.9\,pp); measure the pattern guard via PDR (C3),
not mean Conf.

%=============================================================================
\section{RQ1: Overall Effectiveness (E1)}
\label{sec:results:rq1}
%=============================================================================

\textbf{Claim (C1--C3 jointly).}
The three-insight loop should raise mean Conf.\ over one-shot and
test-feedback baselines; a lower Strict rate under M would support C3
(conjunctive gate), not weaker code.

\subsection{Main Results}

The headline question for RQ1 is whether specification-guided repair improves
\emph{conformance}, and whether a lower strict-success rate is a bug or a
feature of the conjunctive gate.
Table~\ref{tab:main-results} answers with all ten evaluated modes on 120 tasks
each (B0--B5, M, A1--A3).
The figures that follow unpack the same claim: where mass sits in the
conformance distribution, how the CDF orders modes, how many LLM calls the
repair loop spends, and how quality trades against cost.

\begin{table}[H]
  \caption{Main comparison results ($n = 120$ tasks per mode).  Best formal conformance
    shown in bold.  Strict = strict success / acceptance rate
    ($\mathrm{Accept}{=}1$); Conf.\ = mean formal conformance
    ($\overline{\mathrm{Conf}}$); Kill = mutation kill rate (checker precision on
    injected mutants, not generator quality---see \S\ref{sec:results:rq1}); Fail.\ = mean
    strict evaluation failures per task (out of 64).
    B3--B5 from \code{run\_ccf\_b\_extended\_v1} (repeat~0).
    \emph{Note the strict-success paradox:} M's Conf.\ is highest (90.4\%) while
    its Strict column (25.0\%) is slightly below B1/B2---because acceptance
    additionally requires a clear pattern screen (Figure~\ref{fig:accept-tree}).}
  \label{tab:main-results}
  \centering
  \thesistable{%
    \footnotesize
    \renewcommand{\arraystretch}{1.12}
    \begin{tabular*}{\linewidth}{@{\extracolsep{\fill}}l r r r r r@{}}
      \toprule
      Mode & Strict & Conf. & Kill & Lat.\ (ms) & Fail. \\
      \midrule
      B0   & 27.5\%          & 27.5\%        & 79.0\%       & 7.5           & 0.00 \\
      B1   & 26.7\%          & 84.2\%        & 62.3\%       & 9{,}564       & 1.76 \\
      B2   & 26.7\%          & 88.0\%        & 61.4\%       & 10{,}769      & 1.85 \\
      B3   & 5.0\%           & 81.9\%        & 51.0\%       & 13{,}210      & 2.33 \\
      B4   & 5.0\%           & 87.1\%        & 49.3\%       & 10{,}758      & 2.49 \\
      B5   & 5.0\%           & 81.1\%        & 50.6\%       & 6{,}212       & 2.33 \\
      \textbf{M}  & 25.0\%   & \textbf{90.4\%} & 60.6\%     & 43{,}659      & 1.98 \\
      A1   & 26.7\%          & 86.5\%        & 61.9\%       & 41{,}689      & 1.80 \\
      A2   & 26.7\%          & \textbf{90.9\%}& 60.1\%     & 12{,}080      & 1.97 \\
      A3   & 26.7\%          & 84.1\%        & 62.7\%       & 9{,}314       & 1.77 \\
      \bottomrule
    \end{tabular*}%
  }
\end{table}

Table~\ref{tab:main-results} states the aggregate means; the figures below
unpack \emph{where} each mode's mass sits, not only the headline percentages.

\begin{figure}[t]
  \centering
  \widefig{figures/mode_distribution_strict_conformance.pdf}
  \caption{\textbf{Claim:} M concentrates mass in the upper conformance band;
    B0 does not.
    \textbf{How to read:} each violin/box is per-task $\mathit{Conf}(c,t)$ over
    120 tasks; compare the right-hand mass of M with B0's low mode.}
  \label{fig:conf-distribution}
\end{figure}

Figure~\ref{fig:conf-distribution} separates difficulty calibration from
LLM quality: B0's bimodal spread marks template failure on hard tasks, while
B1--M cluster in the 80--100\% band with M shifted furthest right.

\begin{figure}[t]
  \centering
  \widefig{figures/formal_conformance_cdf.pdf}
  \caption{\textbf{Claim:} M and A2 dominate the upper CDF tail.
    \textbf{How to read:} curves farther right mean higher conformance for more
    tasks; B0's steep early rise marks template failure on hard tasks.}
  \label{fig:conf-cdf}
\end{figure}

The CDF in Figure~\ref{fig:conf-cdf} makes the ordering explicit: M and A2
occupy the upper tail, B2 trails slightly, and B1 sits below both repair-heavy
modes---with only 6/120 strict-conformance wins over B2, mean aggregation
and overlap tertiles (E3) are the informative signals.

\begin{figure}[t]
  \centering
  \widefig{figures/summary_metrics_by_mode.pdf}
  \caption{\textbf{Claim:} highest mean Conf.\ coincides with the most LLM calls.
    \textbf{How to read:} left axis = Strict / Conf.; right axis = mean LLM
    calls; M's 2.08 calls show the repair loop is exercised, not idle.}
  \label{fig:summary-metrics}
\end{figure}

Figure~\ref{fig:summary-metrics} links quality to cost at a glance: the modes
with the highest mean conformance (M, A2) are also those that invoke the LLM
more than once per task on average.

\begin{figure}[t]
  \centering
  \widefig{figures/llm_calls_distribution.pdf}
  \caption{\textbf{Claim:} budget $K=3$ is actively used, not decorative.
    \textbf{How to read:} M and A1 mass at 2--3 calls; B1 stays at one.}
  \label{fig:llm-calls}
\end{figure}

The call-count histogram (Figure~\ref{fig:llm-calls}) confirms that budget
$K{=}3$ is exercised: M and A1 concentrate at two or three calls, whereas B1
remains at a single generation attempt.

\begin{figure}[t]
  \centering
  \widefig{figures/radar_metrics.pdf}
  \caption{\textbf{Claim:} M trades latency for conformance; B2 is the cheaper
    moderate alternative.
    \textbf{How to read:} axes are normalised Conf., Strict, kill rate, and
    inverse latency; larger area is not uniformly better.}
  \label{fig:radar-metrics}
\end{figure}

Figure~\ref{fig:radar-metrics} summarises the multi-metric trade-off: no mode
dominates every axis, but M leads on conformance while B2 offers a cheaper
moderate-quality alternative.

Table~\ref{tab:main-results} is the aggregate: M at 90.4\% Conf.\
(+6.2\,pp over B1).
Against B2 (88.0\%), M wins only 6/120 paired tasks on strict conformance
($\approx$5\%; 95\% CI 1--9\%); 110 tasks tie.
The aggregate conformance gap (+2.4\,pp) is \emph{not significant} after Holm correction (Holm $p{=}0.811$, Cliff's
$\delta{\approx}0$).
That is the honest primary-quality reading: structured specification feedback
matters most where overlap is dense (E3, E10), not as a universal headline over
test-feedback.
Wilcoxon on per-task strict success (M vs.\ B1/B2) remains negligible
(Holm-corrected $p{=}1.000$); conformance discriminates M $>$ B2 $>$ B1, with
M vs.\ B1 at Holm $p{=}0.075$ and Cliff's $\delta{=}{+}0.036$ (small).
Extended Self-Refine / Self-Debug / Reflexion-style baselines (B3--B5) sit
below M on conformance (81.1--87.1\%) despite equal $K$ and token budget;
B4 (87.1\%) is closest to B2 (88.0\%) but still trails M by 3.3\,pp.
B3--B5 strict-success rates (5.0\%) reflect the conjunctive gate without
formal witnesses rather than absence of repair iterations.

\subsection{Reading the Distributions}

Three observations settle the rest of RQ1.
First, B0 at 27.5\% is a difficulty calibration: Figure~\ref{fig:conf-distribution}
shows a bimodal template failure, while LLM modes concentrate in the
80--100\% band---headroom the benchmark was built to provide.
Second, the progression B0$\to$B1$\to$B2$\to$M (27.5$\to$84.2$\to$88.0$\to$90.4\%)
is the cost of successively better evidence on \emph{mean} conformance; only
6/120 tasks show strict-conformance wins over B2.
Third, M's \emph{lower} strict success (25.0\% vs.\ 26.7\%) is the conjunctive
gate of Equation~\ref{eq:accept}, not weaker code: acceptance additionally
requires a clear pattern screen, so the accepted set is cleaner.
Mutation kill rate measures checker \emph{precision} against injected syntactic
faults, not generator quality.
B0's 79.0\% exceeds LLM modes (60--63\%) because the template renderer
produces exactly one branch per scenario with no deviation, so every mutant
creates a detectable mismatch; LLM-generated candidates vary in branch order,
variable naming, and formatting, some of which mask injected mutations from
the checker's kill criterion ($\mathrm{Conf}<1$).
Inter-LLM gaps ($\le$2.6\,pp) are small relative to the measurement baseline
and should not be over-read for comparative ranking.

\paragraph{Answer to RQ1 / C1--C3.}
Conformance discriminates M over B1 (+6.2\,pp; marginal Holm $p{=}0.075$).
M beats B2 on 6/120 tasks but the aggregate M--B2 gap is not significant
(Holm $p{=}0.811$); overlap-intensive specifications are where structured
feedback pays (E3, E10).
The slightly lower Strict rate is C3 selectivity, not weaker code.

%=============================================================================
\section{RQ2: Fault Class Prevention (E2, E7)}
\label{sec:results:rq2}
%=============================================================================

\textbf{Claim (C3, with C1).}
The conjunctive gate plus region witnesses should catch ordering/guard
faults that unit-test feedback misses; residual classes name the limits.

\begin{figure}[t]
  \centering
  \widefig{figures/prevention_heatmap_by_mode_eval.pdf}
  \caption{\textbf{Claim:} prevention is complementary across modes, not
    uniformly dominated by one pipeline.
    \textbf{How to read:} rows are modes, columns are mutant operators; darker
    cells mean higher PDR.  Compare M vs.\ B2 within each column.}
  \label{fig:prevention-heatmap}
\end{figure}

Figure~\ref{fig:prevention-heatmap} is the operator-level view: no single mode
wins every column, but M leads on specification-structural faults while B2
remains competitive on boundary-shift operators.

\begin{figure}[t]
  \centering
  \widefig{figures/prevention_bar.pdf}
  \caption{\textbf{Claim:} aggregate PDR alone understates M's structural
    prevention role.
    \textbf{How to read:} bars are aggregate PDR by mode on the full E2
    evaluation (\texttt{prevention\_full\_v1}, $n{=}1{,}704$ mutant-mode
    evaluations per mode); pair with Figure~\ref{fig:prevention-heatmap} for
    per-operator detail.}
  \label{fig:prevention-bar}
\end{figure}

The aggregate bars (Figure~\ref{fig:prevention-bar}) collapse that heterogeneity:
M reaches the highest overall PDR, but the heatmap is needed to see \emph{which}
fault families drive the gap.

\begin{figure}[t]
  \centering
  \widefig{figures/pattern_guard_f1_bar.pdf}
  \caption{\textbf{Claim:} the pattern guard is strongest on ordering and
    guard-inversion faults.
    \textbf{How to read:} bars are F1 by fault category on the E2 mutant
    population; the dashed line marks F1~$=0.6$.}
  \label{fig:pattern-f1}
\end{figure}

\paragraph{Vignette (mutant rejection).}
Take a correct \code{classify\_signal} implementation and apply an SNO mutant
that swaps $S_1$ and $S_2$ in the oracle.
Mode~M's witness $(-3,-10)$ still expects \texttt{Error} under the mutated
order only if the candidate follows the new text; a candidate that kept the
original branch order fails $\mathrm{Accept}$ and counts toward PDR.
A WRO mutant that merely reorders \texttt{if} branches in the candidate is
likewise rejected when the witness set still encodes first-match priority.

What Figure~\ref{fig:prevention-heatmap} adds is complementarity, not a single
winner.
M leads on specification-structural operators (scenario reordering, guard
inversion); B2 leads on boundary-shift operators.
That split is mechanistic: SMT witnesses sit on first-match boundaries and
expose ordering faults, while B2's tests sometimes hit boundary inputs without
naming the violated scenario.

\begin{table}[H]
  \caption{Implementation-screening prevention (E2 full evaluation,
    \texttt{prevention\_full\_v1}, $n{=}852$ mutant evaluations per mode).
    PDR = fraction of faulty mutants correctly rejected; FAR = fraction of
    conforming mutants incorrectly accepted.
    Spec-confusion breakdown appears in Figure~\ref{fig:prevention-heatmap}.}
  \label{tab:prevention}
  \centering
  \thesistable{%
    \begin{tabular}{lcc}
      \toprule
      Mode & PDR & FAR \\
      \midrule
      B1 & 0.0\%   & 100.0\% \\
      B2 & 91.2\%  & 8.8\% \\
      \textbf{M}  & \textbf{95.0\%} & \textbf{5.0\%} \\
      \bottomrule
    \end{tabular}%
  }
\end{table}

Table~\ref{tab:prevention} summarises the implementation-screening track:
B1 lacks an explicit rejection mechanism (FAR 100\%), while M improves both
PDR (+3.8\,pp over B2) and FAR ($-$3.8\,pp) on the full mutant population.
Per-operator complementarity---B2 leading on some boundary-shift operators,
M on specification-structural faults---is visible in
Figure~\ref{fig:prevention-heatmap}, not in the aggregate alone.

Per-category F1 in Figure~\ref{fig:pattern-f1} reinforces the structural
prevention story: Ordering and GuardInversion achieve the highest recall,
confirming that the pattern guard behaves as a specification-derived defect
classifier rather than a loose heuristic rule set.

\paragraph{Answer to RQ2 / C3.}
Ordering and guard-inversion defects are prevented most effectively;
boundary shifts favour B2; arithmetic/output residues remain
(Chapter~\ref{ch:discussion}).

%=============================================================================
\section{RQ5: Quality--Overhead Trade-off (E1)}
\label{sec:results:rq5}
%=============================================================================

\textbf{Claim (practical bound).}
The three insights should buy Conf.\ at a measurable latency premium; a
near-Pareto alternative should exist when the C3 pattern screen is optional.

\begin{figure}[t]
  \centering
  \widefig{figures/mode_distribution_latency.pdf}
  \caption{\textbf{Claim:} repair modes pay a clear latency premium.
    \textbf{How to read:} log-scale wall-clock latency per task; compare the
    M/A1 mass at higher values with B1/B2 near one call.
    B0 and B1 are tightly concentrated near their means.  M and
    A1 exhibit heavier right tails due to repair loop invocations that reach the
    $K=3$ budget limit.}
  \label{fig:latency-distribution}
\end{figure}

The latency distribution (Figure~\ref{fig:latency-distribution}) shows repair
modes paying a clear premium: M and A1 spread toward higher wall-clock times,
while B1/B2 remain near one-call latencies.

\begin{figure}[t]
  \centering
  \widefig{figures/pareto_quality_vs_latency.pdf}
  \caption{\textbf{Claim:} M sits on the quality-dominant side of the
    latency--conformance frontier.
    \textbf{How to read:} right/up is better quality, left is cheaper; the
    dashed polyline is the Pareto frontier and the annotation marks mode~M.}
  \label{fig:pareto}
\end{figure}

The insight in Figure~\ref{fig:pareto} is a trade-off, not a free lunch.
M pays $4.6\times$ vs.\ B2 (43.7\,s vs.\ 10.8\,s), mostly extra LLM calls
(2.08 vs.\ 1.24)---acceptable for parallel CI, not for interactive IDE use.
A2 (12.1\,s, 90.9\% Conf.) is the near-Pareto alternative when pattern
screening is optional; A1 stays slow because weak feedback delays exit.
Mann-Whitney confirms the latency gap is not a sampling artefact
($p{<}0.001$).

\paragraph{Answer to RQ5.}
Best quality at $\approx$4.6$\times$ B2 latency; use A2 when latency binds
and the pattern screen is optional.

%=============================================================================
\section{RQ4 Extended: Boundary Density (E4)}
\label{sec:results:e4}
%=============================================================================

Test-feedback looks strong until witness sampling is held constant.
If denser random tests already covered first-match regions, SMT witnesses
would be optional; Figure~\ref{fig:boundary} demonstrates they are not.

\begin{figure}[t]
  \centering
  \widefig{figures/boundary_coverage_by_density.pdf}
  \caption{\textbf{Claim:} random sampling collapses in the Dense tier; SMT
    does not.
    \textbf{How to read:} coverage vs.\ case budget by density; compare SMT
    vs.\ random within the Dense row at budget~$\ge 8$.}
  \label{fig:boundary}
\end{figure}

At budget 16, random sampling covers far less of the Dense tier than
SMT-guided generation; the gap shrinks in the Sparse tier.
That is why test-feedback baselines look adequate on simple specifications
and degrade on hard ones---the same story as the overlap-tertile split in
Section~\ref{sec:results:rq4}, now measured at the sampling layer alone.

%=============================================================================
\section{Generalisation Study (E8)}
\label{ch:generalization}
%=============================================================================

External validity asks whether SpecIR adapters carry the prevention loop
beyond SOFL syntax, or whether the gains are notation-specific.
Table~\ref{tab:generalisation} and Figure~\ref{fig:generalisation} answer
that question.

\begin{figure}[t]
  \centering
  \widefig{figures/generalisation_bar.pdf}
  \caption{\textbf{Claim:} structured feedback helps across notations, but
    adapter transfer is partial on the expanded $n{=}30$ corpus.
    \textbf{How to read:} bars are mean Conf.\ by notation and mode; compare
    M vs.\ B1 within each group (M leads on both adapters), then M vs.\ B2.}
  \label{fig:generalisation}
\end{figure}

On SOFL/FSF, M leads at 90.4\%.
On the expanded adapter corpora ($n{=}30$ each, \texttt{run\_e8b\_expanded\_v1}),
M beats B1 on Mini-StateMachine (41.6\% vs.\ 24.7\%) and Mini-Z (23.3\% vs.\
5.8\%); M is within 0.8~pp of B2 on Mini-StateMachine and leads B2 by 5.0~pp
on Mini-Z.
Absolute conformance on adapter notations remains low (23--42\%), reflecting
notation-specific generation difficulty rather than SOFL-level overlap structure;
HumanEval/MBPP rows in Table~\ref{tab:generalisation} show stronger transfer on
real-derived guards.
Chapter~\ref{ch:discussion} interprets these boundaries.

\begin{table}[H]
  \caption{\textbf{Key observation:} M leads B1 on all notations; adapter
    absolute conformance is low on the expanded $n{=}30$ corpus (E8b).
    HumanEval-FSF and MBPP-FSF rows report full B1/B2/M from
    \texttt{run\_e8c\_full\_v1}.}
  \label{tab:generalisation}
  \centering
  \thesistable{%
    \small
    \begin{tabular}{lccc}
      \toprule
      Notation & B1 Conf.\% & B2 Conf.\% & M Conf.\% \\
      \midrule
      SOFL/FSF (120 tasks) & 84.2 & 88.0 & \textbf{90.4} \\
      Mini-StateMachine (30 tasks) & 24.7 & 40.8 & \textbf{41.6} \\
      Mini-Z (30 tasks) & 5.8 & 18.3 & \textbf{23.3} \\
      HumanEval-FSF (20 tasks) & 89.7 & \textbf{98.8} & 87.1 \\
      MBPP-FSF (20 tasks) & 90.0 & 95.0 & 95.0 \\
      \bottomrule
    \end{tabular}%
  }
\end{table}

%=============================================================================
\subsection{External Validity: Real-Derived Benchmark (E8c)}
\label{subsec:results:real-derived}
%=============================================================================

Table~\ref{tab:benchmark-by-source} reports conformance across benchmark subsets.
Mode~M reaches 95.0\% on MBPP-FSF (matching B2) and 87.1\% on HumanEval-FSF
(where B2 leads at 98.8\%).
MBPP confirms partial transfer to real-world guard structure; HumanEval shows
that test-feedback baselines can dominate when branches are sparse.

\begin{inlinetable}{Conformance by benchmark source (mode M, $K=3$).}
  \label{tab:benchmark-by-source}
  \footnotesize
  \begin{tabular}{lrrr}
    \toprule
    Subset & $n$ & M Conf (\%) & M Strict (\%) \\
    \midrule
    Synthetic hard (E1)          & 120 & 90.4 & 25.0 \\
    HumanEval-FSF (E8c)          & 20  & 87.1 & 70.0 \\
    MBPP-FSF (E8c)               & 20  & 95.0 & 85.0 \\
    Mini-Z (E8b expanded)        & 30  & 23.3 & --- \\
    Mini-StateMachine (E8b)      & 30  & 41.6 & --- \\
    \bottomrule
  \end{tabular}
\end{inlinetable}

The higher strict-success rate on real-derived tasks (70--85\%) versus
synthetic (25\%) reflects the simpler guard structure of HumanEval/MBPP
problems: real benchmark tasks typically have 3--5 non-overlapping branches,
while synthetic hard tasks are designed with maximal overlap density that
challenges first-match precedence reasoning.
Full B1/B2/M baselines on HumanEval-FSF and MBPP-FSF are reported in
Table~\ref{tab:generalisation} from \texttt{run\_e8c\_full\_v1}.

%=============================================================================
\section{Random Benchmark (E10)}
\label{sec:results:e10}
%=============================================================================

E10 tests whether M's advantage over B2 survives without the Z3 overlap filter
(Section~\ref{sec:exp:e10}).
Table~\ref{tab:e10-random} reports results from
\texttt{run\_e10\_random\_v1} ($n{\approx}100$ per mode after deduplication;
B1 $n{=}102$, B2 $n{=}119$, M $n{=}100$).
On this \emph{unfiltered} sample, B2 leads on mean conformance (88.0\%) and
M (80.5\%) trails---confirming that the M--B2 advantage is \emph{not} universal
off the overlap-heavy hard set.

\begin{table}[H]
  \caption{Random benchmark results (E10).
    Uniformly sampled tasks without Z3 overlap filter; modes B1, B2, M;
    same $K{=}3$ and model as E1.
    M trails B2 on aggregate conformance---advantage is overlap-scoped (E3).}
  \label{tab:e10-random}
  \centering
  \thesistable{%
    \small
    \begin{tabular}{lcccc}
      \toprule
      Mode & Strict (\%) & Conf.\ (\%) & Lat.\ (ms) & Win vs.\ B2 \\
      \midrule
      B1 & 5.9 & 78.1 & 7{,}202 & --- \\
      B2 & 5.0 & 88.0 & 7{,}587 & ref. \\
      M  & 5.0 & 80.5 & 15{,}517 & 1/100 \\
      \bottomrule
    \end{tabular}%
  }
\end{table}

The decision criterion is qualitative: M should match B2 at low overlap and
exceed B2 at high overlap (replicating E3 on an unbiased sample).
If the aggregate gap vanishes, the honest claim narrows to
overlap-intensive specifications only.

\begin{figure}[t]
  \centering
  \widefig{figures/performance_vs_overlap.pdf}
  \caption{E10: mean conformance by overlap-rate tertile on the unfiltered random
    benchmark.  Aggregate means on the full sample favour B2; tertile breakdown
    shows where M gains concentrate (if any).}
  \label{fig:e10-overlap}
\end{figure}

%=============================================================================
\section{External SOFL Benchmark (E11)}
\label{sec:results:e11}
%=============================================================================

E11 evaluates curated, non-synthetic FSF tasks with citable external provenance
(Section~\ref{sec:exp:e11}).
Table~\ref{tab:e11-external} reports repeat~0 results from
\texttt{run\_e11\_external\_v1} ($n{=}22$; zero ID overlap with E1).

\begin{table}[H]
  \caption{External SOFL benchmark (E11).  Hand-curated from textbook and IET
    pattern sources; not generated by \texttt{HardTaskGenerator}.}
  \label{tab:e11-external}
  \centering
  \thesistable{%
    \small
    \begin{tabular}{lcccc}
      \toprule
      Mode & $n$ & Strict (\%) & Conf.\ (\%) & Lat.\ (ms) \\
      \midrule
      B1 & 22 & 63.6 & 89.8 & 2{,}018 \\
      B2 & 22 & 63.6 & 89.8 & 4{,}171 \\
      M  & 22 & 59.1 & 81.9 & 5{,}103 \\
      \bottomrule
    \end{tabular}%
  }
\end{table}

On the external corpus, B1 and B2 tie at 89.8\% mean conformance; M trails
(81.9\%)---these curated tasks have moderate guard overlap and simpler branch
structure than the E1 hard set, consistent with the deployment boundary in
Chapter~\ref{ch:discussion}.
Provenance for each task is recorded in \texttt{benchmarks/external\_sofl.json}.

%=============================================================================
\section{Multi-Seed Stability (E12)}
\label{sec:results:e12}
%=============================================================================

E12 addresses single-seed LLM stochasticity on a 30-task stratified subset
(Section~\ref{sec:exp:e12}).
Table~\ref{tab:e12-stability} reports E12 from
\texttt{run\_e12\_stability\_v1} (30 stratified tasks, seeds $\{0,1,2\}$).
B2 leads on mean conformance (89.2\% vs.\ 86.3\%); M wins 0/30 paired tasks.
Friedman test across seeds: $p{=}0.221$; ranking is \emph{not} seed-stable
(B2 ranks first on seeds 0--1, tie on seed~2).

\begin{table}[H]
  \caption{Multi-seed stability (E12).
    30 stratified tasks, seeds $\{0,1,2\}$, modes B2 and M.
    Win rate and Friedman test assess whether the M--B2 ranking is seed-stable.}
  \label{tab:e12-stability}
  \centering
  \thesistable{%
    \small
    \begin{tabular}{lccc}
      \toprule
      Mode & Mean Conf.\ (\%) & Std.\ across seeds & Ranking vs.\ B2 \\
      \midrule
      B2 & 89.2 & 0.0 & 1st (seeds 0--1) \\
      M  & 86.3 & 16.4 & 2nd / tie \\
      \addlinespace
      \multicolumn{4}{l}{\textit{Paired win rate M vs.\ B2: 0 / 30 (0\%)}} \\
      \multicolumn{4}{l}{\textit{Friedman $p{=}0.221$; ranking stable: No}} \\
      \bottomrule
    \end{tabular}%
  }
\end{table}

%=============================================================================
\section{VerifierLoop-FSF Baseline (E13 / B6)}
\label{sec:results:b6}
%=============================================================================

REV-7 implements mode~B6: SMT witnesses during repair with structured
counterexample prompts but without Semantic Feedback IR or pattern guard
(Section~\ref{sec:exp:methods}).
Table~\ref{tab:b6-stratified} compares B6 with B2 and M on the same 30-task
stratified subset as E12 (\texttt{run\_b6\_stratified\_v1}, repeat~0).
B6 (86.0\% mean conformance) matches B2 (86.2\%) within 0.2~pp while M trails
(83.3\%) on this subset---showing that witness-guided repair without the full IR
recovers test-feedback-level aggregate performance, but does not close the gap to
M on the full hard benchmark (Table~\ref{tab:b6-full}).

\begin{table}[H]
  \caption{VerifierLoop-FSF (B6) vs.\ B2/M on 30-task stratified subset (E13).
    Same tasks and $K{=}3$ as E12; single repeat.}
  \label{tab:b6-stratified}
  \centering
  \thesistable{%
    \small
    \begin{tabular}{lccc}
      \toprule
      Mode & $n$ & Conf.\ (\%) & Lat.\ (ms) \\
      \midrule
      B2 & 30 & 86.2 & 5{,}416 \\
      B6 & 30 & 86.0 & 16{,}245 \\
      M  & 30 & 83.3 & 15{,}099 \\
      \bottomrule
    \end{tabular}%
  }
\end{table}

\begin{inlinetable}{B6 full hard-benchmark run (\texttt{run\_b6\_full\_v1}, $n{=}100$).}
  \label{tab:b6-full}
  \footnotesize
  \begin{tabular}{lrrrr}
    \toprule
    Mode & Strict (\%) & Conf.\ (\%) & Kill (\%) & Lat.\ (ms) \\
    \midrule
    B6 & 5.0 & 78.6 & 52.1 & 13{,}236 \\
    \bottomrule
  \end{tabular}
\end{inlinetable}

On the full 100-task export, B6 reaches 78.6\% conformance versus M's 90.4\% on
the 120-task E1 matrix---the Semantic Feedback IR and pattern guard contribute
the remaining margin beyond raw SMT witnesses.
B6 latency (13.2\,s) sits between B4 (10.8\,s) and B3 (13.2\,s), consistent with
formal checking inside the loop without pattern screening.

%=============================================================================
\section{Sensitivity Analysis}
\label{sec:results:sensitivity}
%=============================================================================

The sensitivity question is whether the reported operating point is brittle.
Figure~\ref{fig:sensitivity-heatmap} maps mean strict formal conformance against
guard-overlap density and repair budget across modes.

\begin{figure}[t]
  \centering
  \widefig{figures/sensitivity_heatmap.pdf}
  \caption{Sensitivity heatmap: mean strict formal conformance as a function of
    task complexity (guard-overlap tertile) and pipeline configuration.  Mode~M
    degrades more gracefully on high-overlap tasks than one-shot baselines.}
  \label{fig:sensitivity-heatmap}
\end{figure}

\subsection{Sensitivity to Task Complexity}

The 120-task benchmark spans a range of complexity in terms of guard-overlap density
and scenario count.
Partitioning tasks into three tertiles based on guard-overlap rate confirms that
higher task complexity degrades conformance across all methods, with M exhibiting
the smallest relative degradation (Section~\ref{sec:results:rq4}).
The sensitivity heatmap in Figure~\ref{fig:sensitivity-heatmap} corroborates this
pattern across all ten primary and ablation modes simultaneously.

\subsection{Sensitivity to Budget $K$}

The budget parameter $K$ controls the maximum number of LLM calls in the repair
loop.
Mode A3 ($K=1$ effectively) achieves 84.1\% conformance; the full mode M ($K=3$)
achieves 90.4\%, an improvement of 6.3\,pp across two additional budget slots.
The marginal return of the second iteration (roughly 3--4\,pp) is higher than that of
the third (roughly 2--3\,pp), consistent with diminishing returns
(Figure~\ref{fig:convergence}).
Extrapolating, $K=2$ would likely yield conformance in the range 87--88\%, comparable
to B2 but with Z3-guided feedback.
The $K=3$ budget appears to be a reasonable operating point: it captures most of the
attainable gain while keeping total LLM calls below 3.0 per task on average.

\subsection{Sensitivity to Formal Case Count}

The formal checker generates up to 64 conformance cases per task for final strict
scoring, but uses only 16 cases per iteration during the repair loop.
The 16-case internal budget trades coverage for latency; the 64-case strict pass
improves sensitivity to guard-overlap faults on the final reported score.
Increasing the internal case count from 16 to 32 would improve counterexample precision
at the cost of approximately 50\% more Z3 time per iteration---a trade-off identified
for future empirical investigation.

\subsection{Sensitivity to LLM Choice}

All experiments used Qwen2.5-27B-Instruct (\textsf{Qwen-27B}).
The experimental design is not specific to this model; the only model-specific
assumptions are (a)~the model can follow a structured Python code generation prompt,
and (b)~the model can incorporate counterexample inputs into a repair prompt.
A stronger base model would likely improve all LLM-based modes
proportionally, with the B1/M gap potentially widening if that model learns to
avoid ordering errors from the first attempt.
Conversely, a weaker model would benefit more from the repair loop, increasing M's
relative advantage.

%=============================================================================
\section{Statistical Test Summary}
\label{sec:results:stats}
%=============================================================================

\begin{table}[t]
\caption{Statistical test summary for E1 (120-task hard benchmark, repeat~0). Wilcoxon signed-rank for per-task strict success and conformance; Mann--Whitney~U for latency. Holm--Bonferroni correction within each metric family. Cliff's $\delta$: positive favours the first mode in each pair.}
\label{tab:stats}
\centering
\small
\begin{tabular}{lccc}
\toprule
Comparison & $p$ (Holm) & Significant ($\alpha{=}0.05$) & Cliff's $\delta$ / ratio \\
\midrule
\multicolumn{4}{l}{\textit{Primary comparisons}} \\
\addlinespace
\texttt{M} vs \texttt{B1} (strict success) & 1.0000 & No & -0.0167 \\
\texttt{M} vs \texttt{B2} (strict success) & 1.0000 & No & -0.0167 \\
\texttt{M} vs \texttt{B1} (conformance) & 0.0745 & No & +0.0358 \\
\texttt{M} vs \texttt{B2} (conformance) & 0.8113 & No & +0.0002 \\
\texttt{M} vs \texttt{B2} (latency) & <0.001 & Yes & 4.05$\times$ \\
\addlinespace
\multicolumn{4}{l}{\textit{All pairwise conformance (Holm-corrected)}} \\
\addlinespace
\texttt{A1} vs \texttt{B0} (Conf) & <0.001 & Yes & +0.4750 \\
\texttt{A1} vs \texttt{B1} (Conf) & 1.0000 & No & +0.0227 \\
\texttt{A1} vs \texttt{B2} (Conf) & 1.0000 & No & -0.0112 \\
\texttt{A1} vs \texttt{M} (Conf) & 1.0000 & No & -0.0119 \\
\texttt{A2} vs \texttt{A1} (Conf) & 0.9931 & No & +0.0285 \\
\texttt{A2} vs \texttt{B0} (Conf) & <0.001 & Yes & +0.5112 \\
\texttt{A2} vs \texttt{B1} (Conf) & 0.0745 & No & +0.0517 \\
\texttt{A2} vs \texttt{B2} (Conf) & 1.0000 & No & +0.0171 \\
\texttt{A2} vs \texttt{M} (Conf) & 1.0000 & No & +0.0174 \\
\texttt{A3} vs \texttt{A1} (Conf) & 1.0000 & No & -0.0272 \\
\texttt{A3} vs \texttt{A2} (Conf) & 1.0000 & No & -0.0568 \\
\texttt{A3} vs \texttt{B0} (Conf) & <0.001 & Yes & +0.4569 \\
\texttt{A3} vs \texttt{B1} (Conf) & 1.0000 & No & -0.0044 \\
\texttt{A3} vs \texttt{B2} (Conf) & 1.0000 & No & -0.0387 \\
\texttt{A3} vs \texttt{M} (Conf) & 1.0000 & No & -0.0407 \\
\texttt{B1} vs \texttt{B0} (Conf) & <0.001 & Yes & +0.4569 \\
\texttt{B2} vs \texttt{B0} (Conf) & <0.001 & Yes & +0.4871 \\
\texttt{B2} vs \texttt{B1} (Conf) & 1.0000 & No & +0.0341 \\
\texttt{M} vs \texttt{B0} (Conf) & <0.001 & Yes & +0.5067 \\
\texttt{M} vs \texttt{B1} (Conf) & 0.0745 & No & +0.0358 \\
\texttt{M} vs \texttt{B2} (Conf) & 0.8113 & No & +0.0002 \\
\addlinespace
\multicolumn{4}{l}{\textit{All pairwise strict success (Holm-corrected)}} \\
\addlinespace
\texttt{A1} vs \texttt{B0} (strict) & 1.0000 & No & -0.0083 \\
\texttt{A1} vs \texttt{B1} (strict) & 1.0000 & No & +0.0000 \\
\texttt{A1} vs \texttt{B2} (strict) & 1.0000 & No & +0.0000 \\
\texttt{A1} vs \texttt{M} (strict) & 1.0000 & No & +0.0167 \\
\texttt{A2} vs \texttt{A1} (strict) & 1.0000 & No & +0.0000 \\
\texttt{A2} vs \texttt{B0} (strict) & 1.0000 & No & -0.0083 \\
\texttt{A2} vs \texttt{B1} (strict) & 1.0000 & No & +0.0000 \\
\texttt{A2} vs \texttt{B2} (strict) & 1.0000 & No & +0.0000 \\
\texttt{A2} vs \texttt{M} (strict) & 1.0000 & No & +0.0167 \\
\texttt{A3} vs \texttt{A1} (strict) & 1.0000 & No & +0.0000 \\
\texttt{A3} vs \texttt{A2} (strict) & 1.0000 & No & +0.0000 \\
\texttt{A3} vs \texttt{B0} (strict) & 1.0000 & No & -0.0083 \\
\texttt{A3} vs \texttt{B1} (strict) & 1.0000 & No & +0.0000 \\
\texttt{A3} vs \texttt{B2} (strict) & 1.0000 & No & +0.0000 \\
\texttt{A3} vs \texttt{M} (strict) & 1.0000 & No & +0.0167 \\
\texttt{B1} vs \texttt{B0} (strict) & 1.0000 & No & -0.0083 \\
\texttt{B2} vs \texttt{B0} (strict) & 1.0000 & No & -0.0083 \\
\texttt{B2} vs \texttt{B1} (strict) & 1.0000 & No & +0.0000 \\
\texttt{M} vs \texttt{B0} (strict) & 1.0000 & No & -0.0250 \\
\texttt{M} vs \texttt{B1} (strict) & 1.0000 & No & -0.0167 \\
\texttt{M} vs \texttt{B2} (strict) & 1.0000 & No & -0.0167 \\
\bottomrule
\end{tabular}
\end{table}


The decision implication of Table~\ref{tab:stats} is metric choice.
Strict-success comparisons among LLM modes are not significant after
Holm correction ($p{=}1.000$): the conjunctive gate is selective by design.
Latency comparisons confirm the expected cost of M over B2
(Mann--Whitney $p{<}0.001$, $4.05\times$ slower).
Conformance gaps (+6.2\,pp over B1) are the adoption signal; the M--B2
aggregate gap is not significant after Holm correction ($p{=}0.811$).
Per-task win rate (6/120) and overlap stratification (E3, E10) supplement
the headline means.
These align with the per-task distributions
in Figures~\ref{fig:conf-distribution} and~\ref{fig:conf-cdf}.

\bigskip

\noindent Taken together: mode~M attains the highest formal conformance
(90.4\%), and the CEGIS repair loop is the largest contributing component.
The latency premium is real but acceptable for CI pre-merge gates.
What remains open---pattern-guard nuance, residual failures, and the
strict-success paradox---needs interpretation rather than more tables.

% -----------------------------------------------------------------------------
\section*{Chapter Summary}
\addcontentsline{toc}{section}{Chapter Summary}
% -----------------------------------------------------------------------------

Four empirical answers organise the evidence.
RQ4/C1: the M--baseline gap grows with overlap density (C1 supported).
RQ3/C2: typed feedback is the largest repair lever (E6 +7.7\,pp; A3 $-$6.3\,pp).
RQ1/C1--C3 jointly: 90.4\% mean Conf.\ (+6.2\,pp over B1; M wins 6/120 vs.\
B2 but Holm $p{=}0.811$ on aggregate conformance).
RQ2/C3: ordering and guard-inversion faults are prevented most effectively;
implementation-screening PDR 95.0\% and FAR 5.0\% confirm that the conjunctive
gate reduces escaped defects.  RQ5: the quality--latency trade-off is
acceptable for CI pre-merge gates.
Chapter~\ref{ch:discussion} interprets what these numbers mean, what still
fails, and how a sprint team should act.
